% Ficha extraída de datos/tabla_prototipos.tex (P-06)
\begin{table}[H]
\centering
\caption{Sexto prototipo - Pantalla de proyección, búsqueda con LLM y voz.}
\label{tab:prototipo-p06}
\small
\begin{tabularx}{\textwidth}{>{\bfseries}l X}
\toprule
ID Prototipo & P-06 \\ \midrule
Nombre & Pantalla de proyección, búsqueda con LLM y voz \\ \midrule
Objetivos & Integrar la pantalla de proyección (ProjectionScreenActivity2): modo resaltado/captura, búsqueda y sincronización papel$\leftrightarrow$PDF, con los botones en la franja superior (zona de mayor validez del eye tracker según Ibaceta, 2023) y la zona de lectura y las cruces en su posición original. \newline Responder los cinco servicios de búsqueda con Gemini 2.5 Flash bajo IWebSearchProvider, en lugar de las APIs vencidas (Wikipedia, Google, Bing, YouTube, Cloud Vision). \newline Integrar lectura en voz alta y comandos de voz (ISpeechProvider) y las páginas Interacción y Búsquedas. \\ \midrule
Descripción & Con este avance el sistema queda completo y equivale al código de <<Proyecto final>> antes de la verificación. Es una extensión respecto del alcance del seminario, donde IWebSearchProvider quedaba solo como interfaz. \\ \midrule
Requisitos Funcionales & B1–B13, C5, C6, C8, C11, C12, RFS1, RFS2 \\ \midrule
Requisitos No Funcionales & RNF-01, RNF-06, RNF-08 \\ \bottomrule
\end{tabularx}
\fuente{Elaboración propia, 2026.}
\end{table}
